\documentclass[11pt,a4paper]{article}
% Preâmbulo comum. Não duplicar pacotes nos módulos.
\usepackage[utf8]{inputenc}
\usepackage[T1]{fontenc}
\usepackage[brazil]{babel}
\usepackage{amsmath,amssymb,amsthm,mathtools}
\usepackage{geometry}
\geometry{a4paper,margin=2.3cm}
\usepackage{microtype}
\usepackage{enumitem}
\usepackage{graphicx}
\usepackage{xcolor}
\usepackage{booktabs}
\usepackage{hyperref}
\hypersetup{colorlinks=true,linkcolor=blue!40!black,urlcolor=blue!50!black}

\setcounter{tocdepth}{2}
\numberwithin{equation}{section}

\theoremstyle{plain}
\newtheorem{teorema}{Teorema}[section]
\newtheorem{proposicao}[teorema]{Proposição}
\newtheorem{lema}[teorema]{Lema}
\newtheorem{corolario}[teorema]{Corolário}

\theoremstyle{definition}
\newtheorem{definicao}[teorema]{Definição}
\newtheorem{exemplo}[teorema]{Exemplo}
\newtheorem{exercicio}{Exercício}[section]

\theoremstyle{remark}
\newtheorem{observacao}[teorema]{Observação}
\newtheorem{contraexemplo}[teorema]{Contraexemplo}
\newtheorem{intuicao}[teorema]{Intuição}

\newcommand{\R}{\mathbb{R}}
\newcommand{\N}{\mathbb{N}}
\newcommand{\Z}{\mathbb{Z}}
\newcommand{\Q}{\mathbb{Q}}
\DeclareRobustCommand{\eps}{\varepsilon}
\DeclareRobustCommand{\vet}[1]{\vec{#1}}
\DeclareRobustCommand{\dif}{\mathrm{d}}
\DeclareRobustCommand{\abs}[1]{\left\lvert #1 \right\rvert}
\DeclareRobustCommand{\norm}[1]{\left\lVert #1 \right\rVert}
\DeclareMathOperator{\grad}{grad}
\DeclareMathOperator{\Div}{div}
\DeclareMathOperator{\rot}{rot}

\newcommand{\base}[1]{\par\medskip\noindent\textbf{Base para relembrar.} #1\par\medskip}
\newcommand{\cuidado}[1]{\par\medskip\noindent\textbf{Quando a hipótese falha.} #1\par\medskip}

\newcommand{\figura}[2]{%
  \begin{figure}[htbp]
  \centering
  \IfFileExists{#1}{%
    \includegraphics[width=0.84\textwidth]{#1}%
  }{%
    \fbox{\parbox{0.8\textwidth}{\centering Figura ainda não gerada:\\ \texttt{#1}}}%
  }
  \caption{#2}
  \end{figure}%
}

\newenvironment{esqueleto}{%
  \par\medskip\noindent\textbf{Esqueleto da demonstração.}\par\nopagebreak
  \begin{enumerate}[leftmargin=1.4em]
}{\end{enumerate}\medskip}

\title{Caderno de prova 2: Rolle, valor médio e teorema fundamental}
\author{}
\date{}
\begin{document}
\maketitle
\tableofcontents

Prove os seis enunciados. No item 2 pode usar o teorema de Weierstrass. Nos itens 5 e 6 pode usar a definição da integral de Riemann para função contínua, a aditividade no intervalo, a comparação e a integral da constante. Cada item pode usar os anteriores. Não há outra indicação.

\section{Fermat interior}

Seja $f\colon D\to\R$ e seja $c$ ponto interior de $D$. Se $f$ tem um extremo local em $c$ e $f$ é derivável em $c$, então $f'(c)=0$.

\vspace{6.5cm}

\section{Rolle}

Seja $f$ contínua em $[a,b]$ e derivável em $(a,b)$, com $a<b$ e $f(a)=f(b)$. Então existe $c\in(a,b)$ tal que $f'(c)=0$.

\vspace{6.5cm}

\section{Valor médio de Lagrange}

Seja $f$ contínua em $[a,b]$ e derivável em $(a,b)$, com $a<b$. Então existe $c\in(a,b)$ tal que
\[
f'(c)=\frac{f(b)-f(a)}{b-a}.
\]
Use a função auxiliar
\[
g(x)=f(x)-f(a)-\frac{f(b)-f(a)}{b-a}(x-a)
\]
e o teorema de Rolle.

\vspace{6.5cm}

\section{Derivada nula}

Seja $I\subseteq\R$ um intervalo e seja $f\colon I\to\R$ contínua em $I$ e derivável no interior de $I$, com $f'(x)=0$ em todo ponto interior. Então $f$ é constante em $I$.

\vspace{6.5cm}

\section{Teorema fundamental, parte 1}

Seja $f$ contínua em $[a,b]$ e defina $F(x)=\int_a^x f(t)\,\dif t$. Prove que $F'(c)=f(c)$ para todo $c\in(a,b)$, que a derivada lateral direita em $a$ é $f(a)$ e que a derivada lateral esquerda em $b$ é $f(b)$.

\vspace{6.5cm}

\section{Teorema fundamental, parte 2}

Seja $f$ contínua em $[a,b]$. Se $F\colon[a,b]\to\R$ satisfaz $F'=f$ em $[a,b]$, com derivada bilateral no aberto e derivadas laterais nos extremos, então
\[
\int_a^b f(x)\,\dif x=F(b)-F(a).
\]

\vspace{6.5cm}

\end{document}
